
We addressed the missing cost-performance analysis in UQ research by framing evaluation as Pareto frontier construction rather than winner-take-all comparison.

Our main contributions are:

\begin{enumerate}
\item Systematic benchmark mapping AUROC versus inference cost for 6 UQ method variants on TruthfulQA selective prediction, establishing that 5 methods are Pareto-optimal across cost zones ($1\times, 3\times, 5\times, 10\times )$.
\end{enumerate}

\begin{enumerate}
\item MC dropout k=3 efficiency point identified (0.704 AUROC at $3\times$ cost). This challenges computer vision conventions (k=30) and establishes k=3 for 8B-scale selective prediction.
\end{enumerate}

\begin{enumerate}
\item Epistemic uncertainty threshold quantified at 8B scale: MC dropout $k\geq 3$ required for AUROC $\geq$ 0.70, while zero-cost methods (temperature scaling 0.682, conformal prediction 0.695) fall short. This indicates post-hoc calibration alone is insufficient when base model has complex miscalibration patterns.
\end{enumerate}

These findings shift UQ evaluation from ''which method wins'' to ''which method for my budget,'' enabling informed trade-offs rather than defaulting to the most expensive option.

\subsubsection{Future Directions}

\textbf{Scale dependency.} Our 8B-only results leave open whether zero-cost methods become competitive at 70B scale. Testing temperature scaling and conformal prediction on Llama-3.1-70B-Instruct would determine if better base calibration eliminates the epistemic uncertainty requirement. If zero-cost methods exceed 0.70 threshold at 70B, practitioners could default to temperature scaling unless high-stakes application requires MC dropout.

\textbf{Adaptive k methods.} MC dropout k=10 showed only +0.006 AUROC versus k=5 despite $2\times$ cost increase. Adaptive k methods with early stopping (when variance converges across forward passes) could reduce average cost while maintaining k=5-equivalent AUROC.

\textbf{Per-query budget allocation.} Five Pareto-optimal methods across cost zones suggest per-query budget allocation. Hybrid approach: use temperature scaling ($1\times$ cost) for easy queries (high confidence), route hard queries (low confidence) to MC dropout k=5 ($5\times$ cost). This dynamic allocation could optimize cost-quality trade-off at query level rather than method level.

\textbf{Cross-dataset transfer.} Conformal prediction AUROC 0.695 suggests in-distribution calibration (TruthfulQA 40\% split) could close the gap and reach threshold. This would establish conformal prediction as zero-cost method viable for $\geq 0.70$ AUROC.

As LLM deployment scales to production applications, budget-aware UQ selection becomes critical infrastructure. Our Pareto frontier framework provides the cost-performance map needed for informed decision-making. Future work at 70B scale and across benchmarks will refine this map, but the core message remains: budget constraints are deployment considerations, and the field should report cost-performance trade-offs rather than winner-take-all rankings.
